%==============================================================================
%  ملحق - الاستغلال والتزامات الخدمة
%==============================================================================

%==============================================================================
\chapter{ملحق: الاستغلال والتزامات الخدمة}
\label{chap:exploitation}
%==============================================================================
النظام المُسلَّم دون التزام خدمة قابل للقياس ليس قابلًا للاستغلال: فمن دون هدف،
لا يمكن وصف أيّ تدهور بأنه حادث. يحدّد هذا الملحق أهداف الخدمة وسياسة النسخ
الاحتياطي وشروط قابلية الرجوع. وهو المرجع لمعايير الاستلام في
الفصل~\ref{chap:delais}.

%------------------------------------------------------------------------------
\section{أهداف مستوى الخدمة (SLO)}
%------------------------------------------------------------------------------
حُدِّدت الأهداف لنشر على \textbf{خادم وحيد}، وفق قرار الاستغلال المعتمَد. وهي
واقعية عن قصد: الوعد بـ 99{,}99~\% دون بنية متكرّرة التزام لا يمكن الوفاء به.

\begin{xltabular}{\textwidth}{>{\bfseries}p{4.2cm} p{3.4cm} X}
\toprule
\textbf{المؤشّر} & \textbf{الهدف} & \textbf{التفصيل} \\
\midrule
\endhead
الإتاحة & 99{,}5~\% شهريًا & نحو 3~س~40~د من التوقّف المقبول شهريًا، خارج الصيانة المُبرمَجة. \\
\addlinespace
الصيانة المُبرمَجة & 4~س شهريًا كحدّ أقصى & يُعلَن عنها قبل 7 أيام، خارج موسم الإنتاج كلّما أمكن. \\
\addlinespace
زمن الاستجابة (p95) & $<$ 500~ms & على شاشات الاطّلاع، عند الحمل الاسمي. \\
\addlinespace
زمن الاستجابة (p99) & $<$ 1~500~ms & يتحمّل الاستعلامات التحليلية الثقيلة. \\
\addlinespace
نسبة الأخطاء & $<$ 1~\% & استجابات 5xx منسوبة إلى مجموع الطلبات. \\
\addlinespace
\textbf{RTO} (زمن الاستئناف) & $<$ 4~س & أقصى مهلة لإعادة الخدمة بعد حادث كبير. \\
\addlinespace
\textbf{RPO} (فقدان المعطيات) & $<$ 1~س & أقصى فقدان مقبول، تضمنه وتيرة النسخ الاحتياطي. \\
\bottomrule
\end{xltabular}

\begin{encadre}
\textbf{الحمل الاسمي المرجعي:} 200 مستخدم متزامن في ذروة الموسم، 10~000 خلية
مزوّدة بحسّاسات، ونحو مليون قياس مُستقبَل يوميًا. لا تصلح أهداف زمن الاستجابة
إلّا ضمن هذه الحدود؛ وما تجاوزها يستوجب مراجعة الأبعاد.
\end{encadre}

%------------------------------------------------------------------------------
\section{النسخ الاحتياطي والاسترجاع}
%------------------------------------------------------------------------------
\subsection{السياسة}
\begin{itemize}
  \item \textbf{نسخ احتياطي متواصل} لسجلّات معاملات PostgreSQL (أرشفة WAL)، بما
        يضمن RPO في حدود ساعة.
  \item \textbf{نسخة كاملة يومية}، مشفّرة عند التخزين.
  \item \textbf{التخزين خارج الموقع} على وسيط منفصل عن خادم الإنتاج --- فالنسخة
        المخزّنة على الآلة التي تحميها لا تحمي شيئًا.
  \item \textbf{الاحتفاظ}: 7 نسخ يومية، و4 أسبوعية، و12 شهرية.
  \item تدخل \textbf{صور الزيارات} وملف \texttt{ConfigZumm.ini} في نطاق النسخ
        الاحتياطي شأنها شأن قاعدة المعطيات.
\end{itemize}

\subsection{اختبار الاسترجاع}
\begin{encadre}
\textbf{النسخة الاحتياطية التي لم تُسترجَع قطّ ليست نسخة احتياطية.} يُختبَر
الاسترجاع الكامل \textbf{شهريًا} على بيئة مخصّصة، ويُؤرشَف التقرير. ونجاح اختبار
الاسترجاع معيار استلام قاطع.
\end{encadre}

يتحقّق الاختبار من ثلاث نقاط: سلامة المعطيات المسترجَعة، واحترام RTO المعلن،
وقدرة المستغلّ على إنجاز الإجراء اعتمادًا على دليل التشغيل وحده.

%------------------------------------------------------------------------------
\section{المراقبة والتنبيهات}
%------------------------------------------------------------------------------
\begin{xltabular}{\textwidth}{>{\bfseries}p{4.2cm} p{3.4cm} X}
\toprule
\textbf{الإشارة} & \textbf{عتبة التنبيه} & \textbf{الإجراء المنتظَر} \\
\midrule
\endhead
تعذّر الوصول إلى التطبيق & إخفاقان متتاليان & تدخّل فوري \\
مساحة القرص & $>$ 80~\% & توسيع أو تنظيف خلال 48~س \\
زمن الاستجابة p95 & $>$ 500~ms خلال 15~د & تحليل الاستعلام المتسبّب \\
نسبة أخطاء 5xx & $>$ 1~\% خلال 5~د & تدخّل فوري \\
إخفاق النسخ الاحتياطي & حالة واحدة & تدخّل في اليوم نفسه \\
شهادة TLS & انتهاء الصلاحية $<$ 30~يومًا & تجديد \\
تأخّر استقبال معطيات الحسّاسات & $>$ 1~س & التحقّق من وسيط MQTT \\
\bottomrule
\end{xltabular}

تستند المراقبة إلى \textbf{OpenTelemetry} وPrometheus وGrafana، وفق الحزمة
التقنية في الملحق~\ref{chap:technologies}.

%------------------------------------------------------------------------------
\section{قابلية الرجوع}
%------------------------------------------------------------------------------
يجب أن يتمكّن العميل من استعادة نظامه دون الاعتماد على المزوّد. تُسلَّم العناصر
التالية عند التسليم، وعند أيّ طلب لاحق:

\begin{itemize}
  \item \textbf{تصدير كامل للمعطيات} بصيغة مفتوحة وموثّقة (SQL وCSV)، متضمّنًا
        الصور والملفات المرفقة؛
  \item \textbf{الشيفرة المصدرية} وتاريخ نسخها؛
  \item \textbf{وثائق الاستغلال} (دليل التشغيل): التشغيل، الإيقاف، الاسترجاع،
        تدوير الشهادات، قراءة التنبيهات؛
  \item \textbf{مخطّطات قاعدة المعطيات} وبرامج الترحيل؛
  \item \textbf{إجراء إعادة النشر} انطلاقًا من مستودع نظيف، مُتحقَّق منه أثناء
        الاستلام.
\end{itemize}

\begin{encadre}
\textbf{التحقّق:} قابلية الرجوع تُبرهَن ولا تُعلَن. تُنجَز إعادة نشر تجريبية من
المستودع على آلة جديدة أثناء الاستلام.
\end{encadre}

%------------------------------------------------------------------------------
\section{نقل الكفاءات}
%------------------------------------------------------------------------------
يومان من نقل الكفاءات مبرمجان عند التسليم، يغطّيان الإدارة اليومية، وقراءة
المراقبة، وإجراء الاسترجاع، والتشخيص من المستوى الأول. ويُستعمل دليل التشغيل
سندًا، ويبقى المرجع بعد النقل.
